\scanpage{10}
on the other hand that $P$ is locally of finite presentation over $k$, it commutes (as afunctor) to \rlap{d}filtering direct limits of rings, and therefore it is known when we know it's restriction to arguments which are of finite type over $k$, and a fortiori \deleted{are} noetherian. For such \tikz[baseline=(s.base)]{\node[inner sep=0pt](s){arguments};\node[anchor=south west,inner sep=1pt](a) at ([xshift=-1em,yshift=7pt]s.north west){\inserted{\handwritten{(i.e. noetherian)}}};\draw ([xshift=-3pt,yshift=-2pt]s.south west) to[out=90,in=180] ([xshift=0pt,yshift=-2pt]a.south west)--([xshift=0pt,yshift=-2pt]a.south east) to[out=0,in=-90] ([xshift=3pt,yshift=3pt]a.south east);}, using the existence theorem for coherent sheaves, we can give a more geometric expression for the right hand side of (18), \deleted{namely}\tikz[baseline=(s.base)]{\node[inner sep=0pt](s){x};} Assume for simplicity $T=\operatorname{Spec}(B)$ affine, and let
\[
C=A\handwritten{\widehat{\mathord{\tikz[baseline=(s.base)]{\node[inner sep=0pt](s){$\mathrm{x}$};\draw[line width=.5pt] ([xshift=.5pt,yshift=1.3pt]s.center) ellipse[x radius=3pt,y radius=4.5pt];}}}}_k B=\lim A_n\mathord{\tikz[baseline=(s.base)]{\node[inner sep=0pt](s){$\mathrm{x}$};\draw[line width=.5pt] ([xshift=.5pt,yshift=1.3pt]s.center) ellipse[x radius=3pt,y radius=4.5pt];}}_k B\ ,
\]
which is an adic noetherian ring augmented to $B$ (the augmentation ideal being an ideal of definition). Consider $\deleted{\unsure{X_{TB}|X_T|Y_B|Y}}\mathord{\tikz[baseline=(s.base)]{\node[inner sep=0pt](s){$=$};\node[anchor=south west,inner sep=1pt](a) at ([xshift=-4em,yshift=7pt]s.north west){\inserted{\handwritten{$X\widehat{\otimes}_k B$}}};\draw ([xshift=-3pt,yshift=-2pt]s.south west) to[out=90,in=180] ([xshift=0pt,yshift=-2pt]a.south west)--([xshift=0pt,yshift=-2pt]a.south east) to[out=0,in=-90] ([xshift=3pt,yshift=3pt]a.south east);}}X\mathord{\tikz[baseline=(s.base)]{\node[inner sep=0pt](s){$\mathrm{x}$};\draw[line width=.5pt] ([xshift=.5pt,yshift=1.3pt]s.center) ellipse[x radius=3pt,y radius=4.5pt];}}_A C$, \tikz[baseline=(s.base)]{\node[inner sep=0pt](s){which};\node[anchor=south west,inner sep=0pt] at ([xshift=-1em,yshift=9pt]s.north west){\inserted{\deleted{$=X\widehat{\mathord{\tikz[baseline=(s.base)]{\node[inner sep=0pt](s){$\mathrm{x}$};\draw[line width=.5pt] ([xshift=.5pt,yshift=1.3pt]s.center) ellipse[x radius=3pt,y radius=4.5pt];}}}_k B$}}};} is a proper scheme over $C$. We then have
\[
\begin{aligned}
(19)\qquad\qquad & \operatorname{Pic}(X\handwritten{\widehat{\mathord{\tikz[baseline=(s.base)]{\node[inner sep=0pt](s){$\mathrm{x}$};\draw[line width=.5pt] ([xshift=.5pt,yshift=1.3pt]s.center) ellipse[x radius=3pt,y radius=4.5pt];}}}}_k B)\handwritten{\xrightarrow{\sim}}\underset{\handwritten{n}}{\underset{\handwritten{\longleftarrow}}{\lim}}\operatorname{Pic}(X_{nB}),\deleted{\unsure{\longrightarrow}P(B)}\handwritten{\operatorname{Pic}(X\widehat{\otimes}_{k}B)/\operatorname{Pic}(B)\hookrightarrow P(B)}
\end{aligned}
\]
and thus \rlap{P}on \uline{noetherian} arguments $B$, $P$ appears as the étale sheaf associated to the functor $B\handwritten{\mapsto}\operatorname{Pic}(X\handwritten{\widehat{\mathord{\tikz[baseline=(s.base)]{\node[inner sep=0pt](s){$\mathrm{x}$};\draw[line width=.5pt] ([xshift=.5pt,yshift=1.3pt]s.center) ellipse[x radius=3pt,y radius=4.5pt];}}}}_k B)\handwritten{/\operatorname{Pic}\unsure{(}B)}$ \tikz[baseline=(s.base)]{\node[inner sep=0pt](s){this};\node[anchor=south west,inner sep=1pt](a) at ([xshift=0em,yshift=7pt]s.north west){\inserted{\handwritten{which is a \uline{separated} presheaf,}}};\draw ([xshift=-3pt,yshift=-2pt]s.south west) to[out=90,in=180] ([xshift=0pt,yshift=-2pt]a.south west)--([xshift=0pt,yshift=-2pt]a.south east) to[out=0,in=-90] ([xshift=3pt,yshift=3pt]a.south east);} statement being complemented by the description given above of the group fcr obstructions for an element of $P(B)$ to \unsure{beleng|belang|belong} to $\operatorname{Pic}(X\handwritten{\widehat{\mathord{\tikz[baseline=(s.base)]{\node[inner sep=0pt](s){$\mathrm{x}$};\draw[line width=.5pt] ([xshift=.5pt,yshift=1.3pt]s.center) ellipse[x radius=3pt,y radius=4.5pt];}}}}_k B)$.

As for $\deleted{Q(B\unsure{)}}B\handwritten{\mapsto}Q(B)$, this appears as the \deleted{\unsure{sehaf|sheaf}} étale sheaf associated with the functor
\[
\begin{aligned}
(20)\qquad\qquad & B\handwritten{\longmapsto}\operatorname{Pic}(X\handwritten{\widehat{\mathord{\tikz[baseline=(s.base)]{\node[inner sep=0pt](s){$\mathrm{x}$};\draw[line width=.5pt] ([xshift=.5pt,yshift=1.3pt]s.center) ellipse[x radius=3pt,y radius=4.5pt];}}}}_k B)/\mathord{\tikz[baseline=(s.base)]{\node[inner sep=0pt](s){$\underline{D}$};\node[anchor=south west,inner sep=3pt](a) at ([xshift=-.2em,yshift=6pt]s.north west){\inserted{\handwritten{$(\operatorname{Pic}(B)+$}}};\draw[line width=.4pt] ([xshift=2pt,yshift=-2pt]a.south west) to[out=160,in=180] ([xshift=-2pt]a.west) to[out=90,in=180] (a.north) to[out=0,in=90] ([xshift=3pt]a.east) to[out=-90,in=0] ([xshift=-2pt,yshift=-2pt]a.south west);\draw ([xshift=-4pt,yshift=-5pt]s.south west) to[out=90,in=-140] ([xshift=-1pt,yshift=-2pt]a.south west);}}(B))\ (\handwritten{\mathrel{\tikz[baseline=(s.base)]{\node[inner sep=0pt](s){$\unsure{\hookrightarrow|\longleftrightarrow}$};\fill ([xshift=1pt,yshift=-5pt]s.south) circle (.3pt);}}}Q(B))\qquad\qquad\mathord{\tikz[x=1em,y=1em,baseline=-.8ex]{\fill (0,0) circle (.3pt);\fill (.13,.19) circle (.3pt);}}
\end{aligned}
\]
where $\underline{D}$ is the twisted constant group over $k$ defined above. For an element of $Q(B)$ to belong to the first hand side there are two successive obstructions, the first is in $H^1(\operatorname{Spec}(B),\underline{D})$, which vanishes \tikz[baseline=(s.base)]{\node[inner sep=0pt](s){if};\node[anchor=south west,inner sep=1pt](a) at ([xshift=-1em,yshift=7pt]s.north west){\inserted{\handwritten{for instance}}};\draw ([xshift=-3pt,yshift=-2pt]s.south west) to[out=90,in=180] ([xshift=0pt,yshift=-2pt]a.south west)--([xshift=0pt,yshift=-2pt]a.south east) to[out=0,in=-90] ([xshift=3pt,yshift=3pt]a.south east);} $\underline{D}$ is constant (for instance $k$ separably closed) and $B$ normal, the second is in the Brauer group of $B$, and vanishes \tikz[baseline=(s.base)]{\node[inner sep=0pt](s){if};\node[anchor=south west,inner sep=1pt](a) at ([xshift=-1em,yshift=7pt]s.north west){\inserted{\handwritten{for instance}}};\draw ([xshift=-3pt,yshift=-2pt]s.south west) to[out=90,in=180] ([xshift=0pt,yshift=-2pt]a.south west)--([xshift=0pt,yshift=-2pt]a.south east) to[out=0,in=-90] ([xshift=3pt,yshift=3pt]a.south east);} $X_o$ \unsure{hes|has} a zero-cycle of degree $1$.

Assume for simplicity that the irreducible components of $X_o$ are
% archive sheet number: 9
